\documentclass[a4paper,11pt]{article}
\usepackage[utf8]{inputenc}
\usepackage[T1]{fontenc}
\usepackage[ngerman]{babel}
\usepackage{amsmath,amssymb}
\usepackage{xcolor}
\usepackage{colortbl}
\usepackage{array}
\usepackage{tikz}
\usepackage{enumitem}
\usepackage{geometry}
\geometry{a4paper,margin=2cm}
\definecolor{bgdark}{HTML}{1E1E1E}
\definecolor{fgtext}{HTML}{E6E6E6}
\definecolor{accent}{HTML}{6CB6FF}
\definecolor{accent2}{HTML}{F2A65A}
\definecolor{gridcol}{HTML}{4A4A4A}
\definecolor{linecol}{HTML}{8A8A8A}
\pagecolor{bgdark}
\color{fgtext}
\arrayrulecolor{linecol}
\newcommand{\karo}[1]{\par\noindent\begin{tikzpicture}\draw[step=5mm, gridcol, very thin] (0,0) grid (\linewidth,#1);\end{tikzpicture}\par}
\newcommand{\aufgabe}[2]{\medskip\par\noindent{\color{accent}\large\textbf{Aufgabe #1}}\hfill{\color{accent2}\small #2}\par\vspace{0.3em}\hrule height0.4pt\vspace{0.6em}}
\newcommand{\kopf}[1]{\noindent{\color{accent}\Large\textbf{Numerik — #1}}\par\vspace{0.4em}\noindent\textbf{Name:}~\rule[-2pt]{5cm}{0.4pt}\hfill\textbf{Datum:}~\rule[-2pt]{3cm}{0.4pt}\par\vspace{0.3em}\hrule height0.8pt\vspace{1em}}
\setlength{\parindent}{0pt}
\begin{document}
\kopf{Fouriertransformation \& FFT}

\textbf{Notation der Vorlesung.} Wellen überlagern sich \emph{linear} (Superposition): die
resultierende Amplitude ist die Summe der Teilamplituden. Gleichphasige Wellen verstärken sich
(konstruktive Interferenz), gegenphasige schwächen sich ab (destruktive Interferenz); die
Einzelwellen selbst bleiben unverändert. Komplexe Welle (Euler):
$A[\cos(\omega t)+i\sin(\omega t)]=Ae^{i\omega t}$. Fourier-Hintransformation (Zeit$\to$Spektrum):
$F(\omega)=\int_{-\infty}^{\infty} f(t)\,e^{-i\omega t}\,dt$, Rücktransformation
$f(t)=\frac{1}{2\pi}\int_{-\infty}^{\infty} F(\omega)\,e^{i\omega t}\,d\omega$. Dabei ist $f(t)$ die
Amplitude über der Zeit und $|F(\omega)|$ der Anteil der Frequenz $\omega$ am Gesamtsignal. Die
\emph{DFT} ist die diskrete Variante für abgetastete Zeitreihen, die \emph{FFT} ein schneller
Algorithmus dafür. Frequenz und Periode hängen über $T=1/f$ zusammen.

\aufgabe{1}{leicht}
An einem festen Ort überlagern sich zwei Sinuswellen gleicher Frequenz:
\[
y_1(t)=3\sin(\omega t),\qquad y_2(t)=5\sin(\omega t).
\]
Beide schwingen gleichphasig. Bestimmen Sie die resultierende Amplitude und geben Sie an, ob
konstruktive oder destruktive Interferenz vorliegt. Wie groß wäre die Amplitude, wenn $y_2$ um
$180^\circ$ phasenverschoben wäre, also $y_2(t)=5\sin(\omega t+\pi)$?
\karo{5cm}

\aufgabe{2}{leicht}
Verwenden Sie die Euler-Formel $Ae^{i\omega t}=A[\cos(\omega t)+i\sin(\omega t)]$ mit Amplitude
$A=2$. Berechnen Sie Real- und Imaginärteil für die Phasen
\[
\omega t=\tfrac{\pi}{2},\qquad \omega t=\pi,\qquad \omega t=\tfrac{\pi}{4}.
\]
Bestätigen Sie insbesondere $e^{i\pi}=-1$.
\karo{5cm}

\aufgabe{3}{leicht--mittel}
Zwei Wellen mit gleicher Frequenz treffen sich:
\[
w_1(t)=4\sin(\omega t),\qquad w_2(t)=4\sin\!\left(\omega t+\tfrac{\pi}{3}\right).
\]
Entscheiden Sie qualitativ, ob sich die Wellen eher verstärken oder abschwächen, und begründen
Sie Ihre Antwort über die Phasenlage. Was passiert im Grenzfall einer Phasendifferenz von $0$,
von $\pi/2$ und von $\pi$?
\karo{5cm}

\aufgabe{4}{mittel}
Erklären Sie, was die Hintransformation $F(\omega)=\int_{-\infty}^{\infty} f(t)\,e^{-i\omega t}\,dt$
inhaltlich abbildet (von welcher Größe in welche?). Worin unterscheidet sich die Bedeutung von
$f(t)$ und von $|F(\omega)|$? Was beschreibt umgekehrt die Rücktransformation?
\karo{5cm}

\aufgabe{5}{mittel}
Eine FFT einer stündlich gemessenen Zeitreihe zeigt Peaks bei den Frequenzen
\[
f_1=0{,}042\ \tfrac{1}{\text{h}},\qquad f_2=0{,}006\ \tfrac{1}{\text{h}},\qquad
f_3=0{,}333\ \tfrac{1}{\text{h}}.
\]
Berechnen Sie zu jedem Peak die Periode $T=1/f$ und interpretieren Sie sie (z.\,B. täglich,
wöchentlich). Welche Zeitskala steckt hinter $f_3$?
\karo{5cm}

\aufgabe{6}{mittel}
Verkehrslärm wird an einer Kreuzung jede Sekunde gemessen. Das FFT-Spektrum zeigt einen sehr
hohen Peak bei einer \emph{niedrigen} Frequenz von $f\approx 1{,}16\cdot 10^{-5}\,$Hz sowie einen
kleineren, breiten Peak bei einer \emph{hohen} Frequenz um $0{,}5\,$Hz. Interpretieren Sie beide
Peaks: Welchem Zyklus entspricht der niederfrequente Peak (Periode berechnen!), und was bedeutet
der hochfrequente Anteil physikalisch?
\karo{6cm}

\aufgabe{7}{mittel}
Warum benötigt man für reale, abgetastete Zeitreihen die diskrete Fouriertransformation (DFT)
statt der kontinuierlichen Integraltransformation $\int_{-\infty}^{\infty} f(t)e^{-i\omega t}dt$?
Welchen Vorteil bringt die FFT gegenüber der direkten Auswertung der DFT-Summe? (Stichwort:
Rechenaufwand.)
\karo{5cm}

\aufgabe{8}{mittel}
Begründen Sie, warum die Fourieranalyse in der Data Science ein nützliches Werkzeug ist, um
Periodizität in Daten zu erkennen. Erläutern Sie am Beispiel des Stromverbrauchs über ein Jahr,
wie sich tägliche, wöchentliche und saisonale Zyklen als Peaks im Spektrum zeigen.
\karo{5cm}

\aufgabe{9}{schwer}
Berechnen Sie von Hand die diskrete Fouriertransformation der Zeitreihe
\[
x=[\,0,\,1,\,0,\,-1\,],\qquad N=4,
\]
nach $X_k=\sum_{n=0}^{N-1} x_n\,e^{-i\,2\pi k n/N}$ für $k=0,1,2,3$. Nutzen Sie
$e^{-i\pi/2}=-i$, $e^{-i\pi}=-1$ usw. Geben Sie die Koeffizienten $X_k$ sowie deren Beträge
$|X_k|$ an und interpretieren Sie, welche Frequenz im Signal steckt.
\karo{7cm}

\end{document}
